\vfill%
\pagebreak

\section*{Exercises}
\markright{\MakeUppercase{Exercises}}%
\subsection*{Exercise 1}

What does the following program print? Work it out by hand, drawing the memory
as the figures of this chapter do, then check your answer by compiling and
running it.

\begin{lstlisting}[style=coloredverbatim]
use std::io;

fn main() {
  let dmut a = copy [1, 2, 3];
  let b = a;
  let dmut c = alias a;
  let dmut d = copy a;
  c[0] = 10;
  d[1] = 20;
  a = copy [7, 8, 9];
  a[2] = 30;
  println(a, " ", b, " ", c, " ", d);
}
\end{lstlisting}

\subsection*{Exercise 2}

The following program should set the three counts to 0, and multiply the two
prices by 3. It contains three errors. Find them, correct them, and compile the
program to test your corrections.

\begin{lstlisting}[style=coloredverbatimError, caption=Ymir program with errors]
use std::io;

fn reset(dmut values: [i32 ; 3]) {
  for i in 0 .. 3 {
    values[i] = 0;
  }
}

fn scale(dmut values: [i32], factor: i32) {
  for mut v in values {
    v = v * factor;
  }
}

fn main() {
  let dmut counts = [4, 5, 6];
  reset(counts);
  let dmut prices = copy [10, 20];
  scale(prices, 3);
  println(counts, " ", prices);
}
\end{lstlisting}

\subsection*{Exercise 3}

Exercise~1 of Chapter~\ref{chap:collections} returned a new slice holding the
elements of another in the reverse order. Write instead a function
\token{reverseInPlace(dmut values: [i32])} that reverses the elements of the
slice it receives, without creating a new one, using the function \token{swap}
of Listing~\ref{lst:mutability:swap}. Check it on a slice of five elements, one
of four, and the empty slice.

\subsection*{Exercise 4}

Write a function \token{removeAll(ref mut values: [i32], unwanted: i32)} that
removes every occurrence of \token{unwanted} from the slice of its caller: after
the call, the variable passed holds only the other elements, in the same order.
Keep a copy of the variable made with \token{let before = readings;} before the
call, and print it after. Does it change? Why?

\subsection*{Exercise 5}

A shop keeps its stock in a map from the name of each item to the number of
items left. Write a function \token{deliver} that adds a delivery, a map of the
same type, to the stock of its caller. An item that the stock does not know yet
is added to it. Test it with a stock of 12 apples and 4 pears, and a delivery of
6 pears and 20 plums.

\subsection*{Exercise 6}

In Listing~\ref{lst:mutability:life}, replace the \token{dcopy} of line~55 with
\token{copy}, then with nothing at all, \token{let before = board;}. Both
programs compile. What do they print, and why?

\subsection*{Exercise 7}

Some shapes of the game of life stop changing after a few generations. Change
\token{step} so that it returns whether the board changed, and write a
\token{main} that steps a row of four living cells on a board of 7 by 7 until it
stops changing, or for 100 generations at most, then prints the generation at
which it stopped and the final board. (\textit{Hint: \token{==} and \token{!=}
  compare two slices of slices element by element.})

\subsection*{Exercise 8}

The \textit{median} of a list of numbers is the number in the middle once the
list is sorted: half of the others are below it, half above. Write a function
\token{median(values: [i32])-> i32} that computes the median of a slice of an
odd length, using the function \token{sort} of Listing~\ref{lst:mutability:sort}.
The slice of the caller must not change. Where does the copy go, and how does
that fit the advice of Section~\ref{sec:mutability:copies}?

\subsection*{Exercise 9}

An interval is a pair of integers, its lowest and its highest value. Write a
function \token{widen(ref dmut bounds: (i32, i32), value: i32)} that widens the
interval of its caller just enough to hold \token{value}, with a \token{match}
whose patterns take references to the fields. Use it to find the lowest and the
highest of a slice of readings, starting from an interval that holds only the
first reading.

\vfill%
\pagebreak

\forceevenpage{}

\section*{Solutions}
\markright{\MakeUppercase{Solutions}}%
\subsection*{Exercise 1}

\begin{lstlisting}[style=bashVerb]
[7, 8, 30] [10, 2, 3] [10, 2, 3] [1, 20, 3]
\end{lstlisting}

\token{b} and \token{c} share the elements of \token{a}, and \token{d} has
elements of its own. \token{c[0] = 10;} writes into the shared elements, which
\token{b} sees too, and \token{d[1] = 20;} changes only \token{d}. Then
\token{a} is given a new slice: from there, \token{a[2] = 30;} changes only the
new elements, and \token{b} and \token{c} keep the first ones.

\subsection*{Exercise 2}

\token{reset} takes an array, which borrows nothing, so a \token{dmut}
parameter is refused: it would change a copy. The function takes a reference
to the array of its caller, \token{ref dmut}, and the call passes
\token{ref counts}. The loop of \token{scale} declares a mutable copy of each
element, which is refused, and a reference iterator cannot go through a
parameter (Section~\ref{sec:mutability:loops}): the loop goes through the
indices. Last, \token{scale} changes the elements of \token{prices}, so the call
passes \token{alias prices}.

\begin{lstlisting}[style=coloredverbatim, escapechar=@, caption=Solution for exercise 2]
use std::io;

fn reset(@\hcb{ref dmut}@ values: [i32 ; 3]) {
  for i in 0 .. 3 {
    values[i] = 0;
  }
}

fn scale(dmut values: [i32], factor: i32) {
  @\hcb{for i in 0 .. values.len \{}@
    @\hcb{values[i] = values[i] * factor;}@
  }
}

fn main() {
  let dmut counts = [4, 5, 6];
  reset(@\hcb{ref}@ counts);
  let dmut prices = copy [10, 20];
  scale(@\hcb{alias}@ prices, 3);
  println(counts, " ", prices);
}
\end{lstlisting}
\vspace{-10pt}%
\begin{lstlisting}[style=bashVerb]
[0, 0, 0] [30, 60]
\end{lstlisting}

\subsection*{Exercise 3}

The loop exchanges the element at index \token{i} with the one at the same
distance from the end, and goes only halfway, since the second half would
exchange the same pairs back. With five elements, the middle one stays in place.

\lstinputlisting[style=coloredverbatim, caption={Solution for exercise 3, \textit{examples/mutability/solutions/reverse.yr}}]{examples/mutability/solutions/reverse.yr}
\vspace{-10pt}%
\begin{lstlisting}[style=bashVerb]
[5, 4, 3, 2, 1]
[4, 3, 2, 1]
[]
\end{lstlisting}

\subsection*{Exercise 4}

The function builds a new slice of the elements it keeps, and gives it to the
variable of its caller through the reference.

\lstinputlisting[style=coloredverbatim, caption={Solution for exercise 4, \textit{examples/mutability/solutions/remove.yr}}]{examples/mutability/solutions/remove.yr}
\vspace{-10pt}%
\begin{lstlisting}[style=bashVerb]
[3, 4, 5]
[3, 0, 4, 0, 0, 5]
\end{lstlisting}

\token{before} does not change. It shares the elements that \token{readings}
held before the call, and \token{removeAll} does not change them: it gives
\token{readings} a new slice, as Section~\ref{sec:mutability:sharing} gave
\token{primes} one.

\subsection*{Exercise 5}

The stock changes in place, so \token{deliver} takes it as a \token{dmut}
parameter, and the call passes it with \token{alias}. The delivery is only read.

\lstinputlisting[style=coloredverbatim, caption={Solution for exercise 5, \textit{examples/mutability/solutions/stock.yr}}]{examples/mutability/solutions/stock.yr}
\vspace{-10pt}%
\begin{lstlisting}[style=bashVerb]
apples: 12
pears: 10
plums: 20
\end{lstlisting}

\subsection*{Exercise 6}

Both programs print the same wrong generations. Generation 1 is:

\begin{lstlisting}[style=bashVerb]
......
..#...
.##...
......
......
\end{lstlisting}

With \token{copy}, \token{before} has a block of rows of its own, but each of
its rows designates the cells of a row of \token{board}
(Figure~\ref{fig:mutability:copies}). With nothing, \token{before} shares
everything with \token{board}. In both cases, when the loop writes the new state
of a cell, \token{before} sees it, and the cells that follow count their
neighbours on a mix of the two generations. Only \token{dcopy} gives
\token{before} cells of its own.

\subsection*{Exercise 7}

\token{step} compares the new board with \token{before} once the loop is done.
The loop of \token{main} counts the generations that changed the board: the
row of four becomes a \textit{beehive}, a shape that no longer changes, at
generation 2.

\lstinputlisting[style=coloredverbatim, caption={Solution for exercise 7, \textit{examples/mutability/solutions/life\_stable.yr}}]{examples/mutability/solutions/life_stable.yr}
\vspace{-10pt}%
\begin{lstlisting}[style=bashVerb]
Stable from generation 2:
.......
.......
..##...
.#..#..
..##...
.......
.......
\end{lstlisting}

\subsection*{Exercise 8}

\token{median} only reads its argument, so its parameter is a plain slice, and
\token{sort} cannot sort it in place: the function copies it, and sorts the
copy. The advice of Section~\ref{sec:mutability:copies} still holds for
\token{sort}, which changes elements and leaves the copy to its caller.
\token{median} is such a caller: it needs the original to stay as it is, so it
makes the copy.

\lstinputlisting[style=coloredverbatim, caption={Solution for exercise 8, \textit{examples/mutability/solutions/median.yr}}]{examples/mutability/solutions/median.yr}
\vspace{-10pt}%
\begin{lstlisting}[style=bashVerb]
Median: 19
Times:  [42, 7, 19, 3, 25]
\end{lstlisting}

\subsection*{Exercise 9}

Each arm takes a reference to the field it may change, and its guard reads the
same field. The interval is a variable of \token{main}, passed as \token{ref
  bounds}; inside \token{widen}, \token{bounds} is a reference, and a variable
that \token{match ref} accepts.

\lstinputlisting[style=coloredverbatim, caption={Solution for exercise 9, \textit{examples/mutability/solutions/bounds.yr}}]{examples/mutability/solutions/bounds.yr}
\vspace{-10pt}%
\begin{lstlisting}[style=bashVerb]
From 6 to 21
\end{lstlisting}
